\chapter{Introduction}
\label{chap:introduction}

This blueprint describes the mathematics formalized in the Lean library
\href{https://github.com/glaroton/LCAnalysis}{LCAnalysis}: the differential and integral
calculus of Michal and Bastiani on locally convex spaces. Every definition and theorem below is
linked to its Lean counterpart, and the dependency graph shows how the results build on each
other. All the statements in this document are proved in Lean, with no \texttt{sorry}, and the
main theorems depend only on the standard axioms of Lean.

\section*{Why Michal--Bastiani calculus}

Infinite-dimensional Lie groups and manifolds of mappings are modeled on spaces that are not
normable: Fréchet spaces of smooth functions such as $C^\infty(M)$, spaces of test functions,
direct limits. On such spaces the Fréchet derivative makes no sense, since there is no norm to
measure the remainder. The calculus that has become standard in infinite-dimensional Lie theory
goes back to Michal \cite{Mic38} and Bastiani \cite{Bas64}, and was developed by Keller
\cite{Kel74}, Hamilton \cite{Ham82}, Milnor \cite{Mil84}, Glöckner \cite{Glo02} and Neeb
\cite{Nee06}. Its main reference today is the book of Glöckner and Neeb \cite{GN26}, which we
follow closely; theorem numbers of the form ``G--N~1.2.13'' refer to it.

The idea is to ask for less than in the Banach case. A map is $\Ck{1}$ if all its directional
derivatives exist and depend continuously on the point \emph{and} on the direction, jointly. No
uniformity in the direction is required. In normed spaces this notion is weaker than Fréchet
differentiability, but the two give the same smooth maps (Chapter~\ref{chap:frechet}).

Mathlib \cite{mC20}, the mathematical library of the Lean proof assistant \cite{dMU21}, has a
complete Fréchet calculus on normed spaces
(\texttt{HasFDerivAt}, \texttt{ContDiff}), and a theory of topological vector spaces and
locally convex spaces. It does not have a differential calculus beyond the normed case. The
purpose of LCAnalysis is to fill this gap, as a foundation for the formalization of
infinite-dimensional Lie theory.

\section*{Standing assumptions}

Unless otherwise stated, $E$, $F$ and $G$ denote real locally convex Hausdorff topological
vector spaces, $U \subseteq E$ is open, and $f \colon U \to F$. The Lean statements sometimes
need less: for instance, the definition of weak integrals makes sense in any topological vector
space. The precise hypotheses can be read in the linked Lean declarations.

Two tools replace norm estimates throughout:
\begin{itemize}
  \item \textbf{Hahn--Banach.} Continuous linear functionals separate points from closed convex
  sets. In particular, a vector is determined by its images under all continuous linear
  functionals.
  \item \textbf{Local convexity.} The closed convex neighborhoods of $0$ form a basis of the
  topology.
\end{itemize}
Most proofs reduce a statement in $F$ to statements about real functions, by composing with
continuous linear functionals; then the one-variable calculus of Mathlib applies.

\section*{What is new and what is not}

The mathematics is classical: every result in this document appears in the cited literature,
and the proofs follow it. The few places where the argument differs from the sources are
pointed out in remarks. As far as we know, Michal--Bastiani calculus had not been formalized
before, neither in Lean nor in another proof assistant.

\section*{How the library was written}

Gabriel Larotonda set the mathematical program, chose the framework and the references, fixed
the design constraints and reviewed the results. The Lean proofs were written by the AI
assistant Claude (Anthropic), working under his direction. Correctness does not rest on trusting
either of them: Lean's kernel checks every proof, and the statements are kept close to the
literature so that they can be compared with it directly.
